% Main-text fragment; supplementary proofs are in supplement_theory.tex.
% Requires amsmath, amssymb, booktabs and natbib.
\section{All-origin loss and computable Bayes risks}
\label{sec:origin-theory}

Direction-only normalization introduces unavoidable classification error throughout a disjoint half-moon family, and introduces additional error in the overlapping configuration. The result concerns every origin in the plane. We first express the lost decisions as cancellation along rays, then identify the geometry that forces this cancellation. A uniform projection limit closes the case of origins escaping to infinity, and signed-density cores make the global lower bound explicit. 

\begin{samepage}
Let $a=0.8$, $b=1.2$, $m_0=(0,0)$ and $m_1=(1,h)$, where $h$ is the upward center displacement. The open class supports are
\begin{align}
 S_0&=\{x:a<\|x\|<b,\ x_2>0\},\nonumber\\
 S_1&=\{x:a<\|x-m_1\|<b,\ x_2<h\}.
 \label{eq:halfmoon-supports}
\end{align}
\end{samepage}
Each support has area $H=\pi(b^2-a^2)/2=2\pi/5$. The primary generator has density $f_i^A=\mathbf1_{S_i}/H$; uniform radius and angle instead gives $f_i^R(x)=\mathbf1_{S_i}(x)/[\pi(b-a)\|x-m_i\|]$. The experiments use equal class priors. The geometric positivity results below hold for arbitrary positive priors $\pi_i$ and densities positive almost everywhere on their respective supports.

\paragraph{Decisions lost along a ray.}
Write $d=\pi_0f_0-\pi_1f_1$, $u_\theta=(\cos\theta,\sin\theta)$ and $[z]_+=\max(z,0)$. Positive $d$ favors class 0 and negative $d$ favors class 1 in the Bayes decision. Along a ray, the factor $r$ is the polar area Jacobian: a fixed angular slice covers more area at greater radius. Its positive and negative signed masses are
\begin{equation}
\begin{aligned}
 A_c(\theta)&=\int_0^\infty[d(c+ru_\theta)]_+r\,dr,\\
 B_c(\theta)&=\int_0^\infty[-d(c+ru_\theta)]_+r\,dr.
\end{aligned}
\label{eq:ray-masses-main}
\end{equation}
Let $R_X^*$ and $R^*_{T_c}$ denote the minimum classification risks from raw coordinates and directions, respectively. Using the masses in~\eqref{eq:ray-masses-main}, the binary Bayes-risk identity yields
\begin{equation}
 \Delta(c):=R^*_{T_c}-R_X^*
 =\int_0^{2\pi}\min\{A_c(\theta),B_c(\theta)\}\,d\theta.
 \label{eq:ray-loss-main}
\end{equation}
Each minimum counts the smaller opposing decision mass merged on that ray. Loss is positive exactly when a positive-measure set of rays combines opposite Bayes preferences. This specializes the established decision-preservation framework \citep{Silva2012,Silva2024} to the origin-dependent direction map. Supplementary Section~\ref{supp:bayes} derives the identity.

\paragraph{Theorem 1: an all-origin transition.}
For $0<h\leq\sqrt{3/5}$, the supports have disjoint interiors and
\begin{equation}
 R_X^*=0,\qquad \Delta(c)>0
 \quad\text{for every }c\in\mathbb R^2.
 \label{eq:all-origin-disjoint}
\end{equation}
At $h=0$, every origin with $c_2=0$ preserves zero risk. More generally, for $h<0$, origins in the separating strip $h\leq c_2\leq0$ preserve zero risk. At $h=1$, the supports overlap and
\begin{equation}
 R_X^*>0,\qquad \Delta(c)>0
 \quad\text{for every }c\in\mathbb R^2.
 \label{eq:all-origin-overlap}
\end{equation}

\emph{Proof sketch.}
For $0\leq h\leq a$, horizontal slicing reduces disjointness to
$\sqrt{a^2-h^2}\geq1/5$, giving the stated threshold. Throughout its positive range, the connected supports intrude into each other's convex hull: $(1,h/2)\in S_0\cap\operatorname{int}\operatorname{conv}(S_1)$ and its reflection belongs to the reverse intersection. The relevant geometric criterion is a probabilistic adaptation of the separability geometry of Hurtado et al.\ \citep{Hurtado2001}. Two connected open regions with mutual convex-hull intrusion have overlapping direction images from every origin. Indeed, disjoint images would be open arcs, at least one of length at most $\pi$; the convex cone generated by that arc would contain its region's convex hull and exclude the other region. Positive densities turn the resulting open angular overlap into positive loss in~\eqref{eq:ray-loss-main}.

For $h=1$, the same argument applies to connected regions exclusive to each class. One class has a connected backbone through $(1,1/4)$, $(-1,1/20)$ and $(0,11/10)$, with a radius-$1/20$ neighborhood contained in its exclusive region. Reflection about $(1/2,1/2)$ produces the other backbone. Strictly positive barycentric coordinates establish mutual hull intrusion. This forces additional loss even though the raw classes already overlap. Supplementary Sections~\ref{supp:geometry}--\ref{supp:overlap} give the complete constructions and proofs.

% Suggested main-text replacement/extension for the old positive-infimum theorem.
% The supplement fragment defines supp:projection-boundary.
\paragraph{Theorem 2: the decision problem at distant origins.}
A positive loss at each fixed origin can still approach zero along an unbounded sequence of origins. To determine a global floor, we also control these escaping sequences. Moving the origin far from the data makes its rays almost parallel. Orthogonal projections are the geometric boundary of radial projections \citep{Orponen2019}; the following result identifies their classification-risk limit. It applies to binary populations with densities in any dimension $n\geq2$. Fix a reference point $p$, let $e$ be a unit vector, and let $P_e$ project onto the hyperplane perpendicular to $e$. Then
\begin{equation}
 \lim_{R\to\infty}\sup_e
 \left|R^*_{T_{p+Re}}-R^*_{P_e}\right|=0.
 \label{eq:projection-limit-main}
\end{equation}
Write $\Delta_\infty(e)=R^*_{P_e}-R_X^*$. The loss surface, completed by these projection limits, is continuous on a compact domain. Consequently,
\begin{equation}
 \inf_c\Delta(c)>0
 \quad\Longleftrightarrow\quad
 \begin{cases}
 \Delta(c)>0&\text{all finite }c,\\
 \Delta_\infty(e)>0&\text{all directions }e.
 \end{cases}
 \label{eq:global-criterion-main}
\end{equation}
Thus optimizing the origin cannot make loss vanish when both finite radial views and all limiting linear projections lose decisions.

\emph{Half-moon consequence.}
The connected winner regions in Theorem~1 have overlapping projection intervals in every direction. Equations~\eqref{eq:projection-limit-main}--\eqref{eq:global-criterion-main} therefore strengthen its conclusion to
\begin{equation}
 \inf_{c\in\mathbb R^2}\Delta(c)>0
 \quad\text{for }0<h\leq\sqrt{3/5}\text{ and for }h=1.
 \label{eq:global-positive}
\end{equation}
This holds for every fixed pair of positive priors and densities positive almost everywhere on their respective supports. The proof controls signed probability mass in total variation, so it establishes convergence of the optimal decision risk. Supplementary Section~\ref{supp:projection-boundary} proves the general result and gives an explicit convergence bound for bounded-support densities of bounded variation.

% Standalone main-text replacement paragraph. Definitions of d and Delta precede it.
\paragraph{A constructive certificate over all origins.}
The loss floor can be bounded using regions where the preferred class has a fixed probability margin. In the plane, let $K_+,K_-$ be compact connected cores whose radius-$\rho$ neighborhoods satisfy $d\geq\kappa_+>0$ and $d\leq-\kappa_-<0$, respectively. Suppose each core meets the other's convex hull, their convex hulls contain a common disk $B(p,\delta)$, and both cores lie within distance $M$ of $p$. Writing $\kappa=\min(\kappa_+,\kappa_-)$ gives
\begin{equation}
 \inf_c\Delta(c)\geq\kappa\min\left\{
 \frac{\sqrt3\rho^3}{6M},\;2\rho\delta\right\}>0.
 \label{eq:constructive-global}
\end{equation}
For bounded origins, mutual hull intrusion supplies a ray through both cores. Their neighborhoods give a common angular interval carrying a fixed amount of opposing signed mass. For exterior origins, rays through the common hull disk supply the bound. Supplementary Section~\ref{supp:constructive-global} proves the result and optimizes the split between these regions. For equal-prior area sampling, optimized whole-plane lower certificates exceed $0.0205$ percentage points (pp) at $h=0.5$ and $0.00119$~pp at $h=1$. Eroded half-annuli provide such cores throughout the disjoint regime of Theorem~1 under both sampling laws and any fixed positive priors.

\paragraph{Fixed-origin risk calculation.}
For the half-annuli, circle tangencies and boundary intersections partition directions into finitely many cells. Radial support order is fixed inside a cell, so one representative ray determines whether opposite preferences coexist. Exact radial primitives then reduce Bayes risk to one angular integral. Ball arithmetic and a bounded-variation truncation bound certify its value. Supplementary Sections~\ref{supp:finite}--\ref{supp:bv} give the finite preservation test, integration formulas and error bounds. At equal priors, the density ratio $f_i^R/f_i^A\in[5/6,5/4]$ also gives $\Delta_R(c)\geq(5/6)\Delta_A(c)$.

\paragraph{One bit repairs the disjoint midpoint representation.}
At $h=0.5$, direction alone has positive loss at every origin. At the midpoint $c=(1/2,1/4)$, append the single radial bit $q=\mathbf1\{\|x-c\|>4/5\}$. For $q=0$, the sign of $9u_1+10u_2$ determines the class; for $q=1$, the sign of $u_2-u_1$ determines it. Positive signs identify class~0. Figure~\ref{fig:bit-repair} shows how the radial bit assigns observations to the two angular rules. Together, the two rules correctly label the complete class supports under their respective bit conditions, giving zero Bayes risk under either sampling law and any positive priors. Supplementary Section~\ref{supp:one-bit} proves the separating inequalities. Since zero additional radial bits leave the positive direction-only floor, one bit is the minimal radial addition at this geometry.

\paragraph{Certified thresholds and radius-error tolerance.}
With the same angular rules, recovery is exact precisely for $\alpha\leq t\leq\beta$, where $\alpha\approx0.7433034374$ and $\beta\approx0.8120914734$ (Threshold Addendum, Proposition~1). If the measured radius is $\widehat r=r+e$, $|e|\leq\varepsilon$, every support label is preserved for every admissible error function exactly when
\begin{equation}
\alpha+\varepsilon\leq t\leq\beta-\varepsilon.
\label{eq:radius-error}
\end{equation}
The midpoint threshold $t_\star=(\alpha+\beta)/2\approx0.7776974554$ maximizes the guaranteed tolerance at $(\beta-\alpha)/2\approx0.0343940180$ (Addendum, Corollary~1). The neural experiments retain $t=0.8$.
